
\documentclass{article}
\usepackage{colortbl}
\usepackage{makecell}
\usepackage{multirow}
\usepackage{supertabular}

\begin{document}

\newcounter{utterance}

\centering \large Interaction Transcript for game `wordle', experiment `medium\_frequency\_words\_no\_clue\_no\_critic', episode 10 with deepseek{-}v4{-}pro.
\vspace{24pt}

{ \footnotesize  \setcounter{utterance}{1}
\setlength{\tabcolsep}{0pt}
\begin{supertabular}{c@{$\;$}|p{.15\linewidth}@{}p{.15\linewidth}p{.15\linewidth}p{.15\linewidth}p{.15\linewidth}p{.15\linewidth}}
    \# & \multicolumn{2}{c}{Player} && \multicolumn{2}{c}{Game Master} \\
    \hline

    \theutterance \stepcounter{utterance}  
    & & & \multicolumn{4}{p{0.6\linewidth}}{
        \cellcolor[rgb]{0.9,0.9,0.9}{
            \makecell[{{p{\linewidth}}}]{
                \texttt{\tiny{[P1$\langle$GM]}}
                \texttt{Vous êtes un magicien du langage qui aime deviner des mots en utilisant les règles données.} \\
\\ 
\texttt{Bienvenue à Wordle ! Vous avez six tentatives pour deviner le mot cible, un mot valide de cinq lettres minuscules. Veuillez utiliser les balises "explication :" et "proposition :" pour fournir une explication concise pour chaque proposition.} \\
\\ 
\texttt{Par exemple, si votre proposition est "année", votre réponse devrait être} \\
\texttt{explication : c'est un mot commun de cinq lettres, et je commence ma proposition avec ce mot.} \\
\texttt{proposition : année} \\
\\ 
\texttt{Après chaque proposition, votre réponse sera validée, et vous recevrez un retour indiquant quelles lettres sont correctes (vert), quelles lettres sont correctes mais à la mauvaise position (jaune), et quelles lettres sont incorrectes (rouge). Ce retour peut être utile pour déterminer quelles lettres inclure ou exclure dans votre prochaine proposition.} \\
\\ 
\texttt{Par exemple, le retour pour "année" pourrait être :} \\
\texttt{retour\_proposition : a<jaune> n<jaune> n<vert> é<jaune> e<rouge>} \\
\\ 
\texttt{L'explication doit contenir des détails sur la manière dont le retour\_proposition est utilisé pour arriver à une nouvelle proposition.} \\
\\ 
\texttt{Ne générez aucun autre texte.} \\
\\ 
\texttt{Commençons avec votre première proposition.} \\
            }
        }
    }
     \\ \\

    \theutterance \stepcounter{utterance}  
    & \multicolumn{4}{p{0.6\linewidth}}{
        \cellcolor[rgb]{0.561,0.737,0.561}{
            \makecell[{{p{\linewidth}}}]{
                \texttt{\tiny{[P1$\rangle$GM]}}
                \texttt{explication : Je commence par un mot commun contenant des voyelles variées et des consonnes fréquentes pour cerner rapidement les lettres potentielles.  } \\
\texttt{proposition : table} \\
            }
        }
    }
    & & \\ \\

    \theutterance \stepcounter{utterance}  
    & & & \multicolumn{2}{p{0.3\linewidth}}{
        \cellcolor[rgb]{0.9,0.9,0.9}{
            \makecell[{{p{\linewidth}}}]{
                \texttt{\tiny{[GM$|$GM]}}
                \texttt{attempts = 1} \\
\texttt{target\_word = revue} \\
\texttt{guess = table} \\
\texttt{guess\_feedback = t<red> a<red> b<red> l<red> e<green>} \\
            }
        }
    }
    & & \\ \\

    \theutterance \stepcounter{utterance}  
    & & & \multicolumn{4}{p{0.6\linewidth}}{
        \cellcolor[rgb]{0.9,0.9,0.9}{
            \makecell[{{p{\linewidth}}}]{
                \texttt{\tiny{[P1$\langle$GM]}}
                \texttt{retour\_proposition : t<rouge> a<rouge> b<rouge> l<rouge> e<vert>} \\
\\ 
\texttt{Fournissez votre réponse uniquement dans ce format.} \\
\texttt{explication : détails} \\
\texttt{proposition : mot} \\
            }
        }
    }
     \\ \\

    \theutterance \stepcounter{utterance}  
    & \multicolumn{4}{p{0.6\linewidth}}{
        \cellcolor[rgb]{0.561,0.737,0.561}{
            \makecell[{{p{\linewidth}}}]{
                \texttt{\tiny{[P1$\rangle$GM]}}
                \texttt{explication : Seul le E est vert en position 5. Les lettres T, A, B, L sont exclues. Je cherche un mot de 5 lettres terminant par E, contenant d'autres voyelles comme I, O, U et des consonnes fréquentes comme R, S, N, C.} \\
\texttt{proposition : crine} \\
            }
        }
    }
    & & \\ \\

    \theutterance \stepcounter{utterance}  
    & & & \multicolumn{2}{p{0.3\linewidth}}{
        \cellcolor[rgb]{0.9,0.9,0.9}{
            \makecell[{{p{\linewidth}}}]{
                \texttt{\tiny{[GM$|$GM]}}
                \texttt{Error: The guessed word is not a valid word for this game.} \\
            }
        }
    }
    & & \\ \\

    \theutterance \stepcounter{utterance}  
    & & & \multicolumn{4}{p{0.6\linewidth}}{
        \cellcolor[rgb]{0.9,0.9,0.9}{
            \makecell[{{p{\linewidth}}}]{
                \texttt{\tiny{[P1$\langle$GM]}}
                \texttt{Votre proposition n'est pas un mot valide pour ce jeu. Veuillez réessayer.} \\
\\ 
\texttt{Fournissez votre réponse uniquement dans ce format.} \\
\texttt{explication : détails} \\
\texttt{proposition : mot} \\
            }
        }
    }
     \\ \\

    \theutterance \stepcounter{utterance}  
    & \multicolumn{4}{p{0.6\linewidth}}{
        \cellcolor[rgb]{0.561,0.737,0.561}{
            \makecell[{{p{\linewidth}}}]{
                \texttt{\tiny{[P1$\rangle$GM]}}
                \texttt{explication : Seul le E est vert en position 5. Les lettres T, A, B, L sont exclues. Je teste un mot valide terminant par E avec des lettres fréquentes non encore essayées comme R, S, N, O.} \\
\texttt{proposition : ronce} \\
            }
        }
    }
    & & \\ \\

    \theutterance \stepcounter{utterance}  
    & & & \multicolumn{2}{p{0.3\linewidth}}{
        \cellcolor[rgb]{0.9,0.9,0.9}{
            \makecell[{{p{\linewidth}}}]{
                \texttt{\tiny{[GM$|$GM]}}
                \texttt{attempts = 2} \\
\texttt{target\_word = revue} \\
\texttt{guess = ronce} \\
\texttt{guess\_feedback = r<green> o<red> n<red> c<red> e<green>} \\
            }
        }
    }
    & & \\ \\

    \theutterance \stepcounter{utterance}  
    & & & \multicolumn{4}{p{0.6\linewidth}}{
        \cellcolor[rgb]{0.9,0.9,0.9}{
            \makecell[{{p{\linewidth}}}]{
                \texttt{\tiny{[P1$\langle$GM]}}
                \texttt{retour\_proposition : r<vert> o<rouge> n<rouge> c<rouge> e<vert>} \\
\\ 
\texttt{Fournissez votre réponse uniquement dans ce format.} \\
\texttt{explication : détails} \\
\texttt{proposition : mot} \\
            }
        }
    }
     \\ \\

    \theutterance \stepcounter{utterance}  
    & \multicolumn{4}{p{0.6\linewidth}}{
        \cellcolor[rgb]{0.561,0.737,0.561}{
            \makecell[{{p{\linewidth}}}]{
                \texttt{\tiny{[P1$\rangle$GM]}}
                \texttt{explication : Le R est vert en position 1 et le E est vert en position 5. Les lettres O, N, C sont exclues, tout comme T, A, B, L des essais précédents. Il me faut un mot en R \_ \_ \_ E, sans O, N, C, T, A, B, L. J'essaie un mot avec des lettres restantes comme S, M, P, D, I, U.} \\
\texttt{proposition : ruide} \\
            }
        }
    }
    & & \\ \\

    \theutterance \stepcounter{utterance}  
    & & & \multicolumn{2}{p{0.3\linewidth}}{
        \cellcolor[rgb]{0.9,0.9,0.9}{
            \makecell[{{p{\linewidth}}}]{
                \texttt{\tiny{[GM$|$GM]}}
                \texttt{Error: The guessed word is not a valid word for this game.} \\
            }
        }
    }
    & & \\ \\

    \theutterance \stepcounter{utterance}  
    & & & \multicolumn{4}{p{0.6\linewidth}}{
        \cellcolor[rgb]{0.9,0.9,0.9}{
            \makecell[{{p{\linewidth}}}]{
                \texttt{\tiny{[P1$\langle$GM]}}
                \texttt{Votre proposition n'est pas un mot valide pour ce jeu. Veuillez réessayer.} \\
\\ 
\texttt{Fournissez votre réponse uniquement dans ce format.} \\
\texttt{explication : détails} \\
\texttt{proposition : mot} \\
            }
        }
    }
     \\ \\

    \theutterance \stepcounter{utterance}  
    & \multicolumn{4}{p{0.6\linewidth}}{
        \cellcolor[rgb]{0.561,0.737,0.561}{
            \makecell[{{p{\linewidth}}}]{
                \texttt{\tiny{[P1$\rangle$GM]}}
                \texttt{explication : R vert en 1, E vert en 5. Lettres exclues : T, A, B, L, O, N, C. Je dois trouver un mot valide en R \_ \_ \_ E avec les lettres disponibles comme S, M, P, D, I, U, G, F, H, V. Je tente un mot courant comme "rime".} \\
\texttt{proposition : rime} \\
            }
        }
    }
    & & \\ \\

    \theutterance \stepcounter{utterance}  
    & & & \multicolumn{2}{p{0.3\linewidth}}{
        \cellcolor[rgb]{0.9,0.9,0.9}{
            \makecell[{{p{\linewidth}}}]{
                \texttt{\tiny{[GM$|$GM]}}
                \texttt{Error: The length of the guessed word is not 5.} \\
            }
        }
    }
    & & \\ \\

    \theutterance \stepcounter{utterance}  
    & & & \multicolumn{2}{p{0.3\linewidth}}{
        \cellcolor[rgb]{0.9,0.9,0.9}{
            \makecell[{{p{\linewidth}}}]{
                \texttt{\tiny{[GM$|$GM]}}
                \texttt{game\_result = ABORT} \\
            }
        }
    }
    & & \\ \\

\end{supertabular}
}

\end{document}
